\subsection{Opened literature and exact scope}


\begin{statement}{P.1}{Known}{The moduli-group interpretation supplies an inclusion}
With inertia preservation and geometric symmetric-group compatibility included,
\citet[Main Theorem and Corollary]{LitHarbaterSchneps2000} identify
the four-point outer automorphism group with $\GZ$ and the corresponding
$n$-point group, $n\ge5$, with $\GT$. The forgetful map supplies
$\GT\hookrightarrow\GZ$.
\end{statement}
\begin{proof}[Derivation and scope]
The four-point pure fundamental group is the marked $\F$ of the prompt.
Its three peripheral classes are $x,y,(xy)^{-1}$, and its geometric
symmetric action is the specified $S_3$. The source's Proposition 4
turns the cyclic compatibility into the three-cycle equation; the
transposition compatibility gives the two-cycle equation. At five
points the remaining compatibility is the pentagon. Restricting an
actual compatible automorphism along the forgetting maps yields the
injection in the source's corollary. This identifies the groups and
maps used here, but the cited theorem does not assert that the injection
is proper. The proof of properness in P.11 below uses the supplied R1
and independently constructed pro-prime inputs.
\end{proof}

\begin{statement}{P.2}{Known}{Categorical coherence and its faithful scope}
Drinfeld's group pentagon is the universal associativity coherence
condition for the modification by a pair $(\lambda,f)$; a single
pointed braided category need not detect that group relation.
\end{statement}
\begin{proof}[Derivation]
In \citet[Section 4, equations (4.5), (4.10), (4.12), (4.13)]{LitDrinfeld1991},
the associator changes by evaluation of $f$ on the two pure monodromies,
and the braiding changes by the multiplier. The free braid situation
retains the group equation. Passing to a category applies its braid
representation to the equation, which is a valid implication in that
direction. The representation may have a kernel. In the category used
in A.2(c), the two relevant monodromies commute, so every $f\in\F'$
evaluates to one; the explicit proof and finite coherence checks below
show that this loss of information occurs here. Thus universal
coherence cannot be inferred back from this one example.
\end{proof}

\begin{statement}{P.3}{Known}{The classical torsion result used here is a centralizer theorem}
For $c=(-1,1)$, \citet[Proposition 4(ii)]{LitLochakSchneps1997}
prove $C_{\GT}(c)=\langle c\rangle$. This result does not itself
classify every torsion element of $\GT$ or $\GZ$. We use the supplied,
refereed R7 for the assertion that every nonidentity finite-order
element of either group is an involution of multiplier $-1$.
\end{statement}
\begin{proof}[Derivation and separation of assertions]
The source's $c$ is the pair with inverted peripheral generators and
trivial $f$, matching our marked normalization. Its centralizer assertion
therefore transfers literally. A centralizer statement at this element
cannot imply that a different torsion element is conjugate to it.
R7 separately gives torsion-freeness of $\HH$, the order-two conclusion,
and $\Delta((-1,f))=\delta((1,f))$. R4 gives zero on every conjugate
of $c$, because $\Delta$ is a character and $c$ is Galois. No general
involution-conjugacy theorem is claimed or needed in this return.
\end{proof}

\begin{statement}{P.4}{Known}{The currently available mapping-class theorems retain comparison hypotheses}
The braid, configuration-space, and higher-genus automorphism theorems
listed below do not supply the CM comparison needed to evaluate
$\Delta|_{\GT}$.
\end{statement}
\begin{proof}[Exact reductions]
\citet[Theorems A and B]{LitMinamideNakamura2022} treat profinite braid
groups on at least four strands and their central quotients. Their
additional central-unit factor disappears only after that specified
quotient; the theorem cannot be specialized to the three-strand modular
source. \citet[Corollaries B and C]{LitHoshiMinamideMochizuki2022}
recover the configuration-space fibre structure and the GT factor in
the outer automorphism group under their stated number-of-points and
prime-set hypotheses. Identification of that factor gives no equality
with Galois or calculation of its abelianization.

\citet[Theorem 1.1]{LitLochakNakamuraSchneps2026} construct full
profinite actions in genus $g>0$ with zero or one marked point.
Their Section 8 retains the required cutting and Lego compatibilities
as further work. Therefore the existence of these individual actions
cannot identify two actions defined by different CM constructions.
The exact low-genus sewing application of
\citet[Theorem 2.3 and Proposition 2.6(i)]{LitBoggi2026v6} is proved
in Part B; it is an identification of the relevant groups with their
maps, and supplies no equality of arbitrarily specified actions.
\citet[Theorem A and Corollary B]{BoggiSymmetric2026} give the
congruence property for finite-subgroup centralizers when the quotient
surface has genus at most two. That result controls the induced
profinite topology, not the residual frame character.
\end{proof}

\begin{statement}{P.5}{Known}{Special-locus results and the exact Galois-only comparison}
The harmonic, equianharmonic and periodic-frame constructions occur
in the opened literature, but none of the specific cited results
identifies the two CM coordinates for every element of $\GT$.
\end{statement}
\begin{proof}[Derivation and quantifiers]
\citet[Theorems 1.1--1.2]{LitNakamuraSchneps2000} define the subgroup
with the extra relations $(\mathrm{III}')$ and $(\mathrm{IV})$ and
prove the Galois image lies in it. The harmonic and equianharmonic
identities are printed in the primary report
\citet[Theorems 1--2]{LitTsunogaiNakamura2002}; the later 2003
article's mathematical full text was not accessible here and is marked
unverified in the bibliography. \citet[Theorem 3.4]{LitSchneps2006}
deduces the further identities for its precisely defined special-locus
subgroup, including $\lambda=1$, $\rho_2=\rho_3=0$ and two additional
compatibilities. Those restrictions cannot be removed by citing the
theorem.

\citet[Proposition 4.4]{LitNakamuraTsunogaiYasuda2010} give an
all-GT triangle action, while Propositions 4.11 and 4.13 and
Theorem 4.16 compare it with the standard action on Galois. The
Galois restriction is explicit; Section 4.12 is not a proposition.
Part B derives its sign with the central powers retained.
\citet[Theorems A and B]{Collas2012GenusZero} give geometricity and
cyclotomic conjugacy for prime-order torsion in genus zero. The
separate genus-one text \citep{CollasGenusOneAuthorText} originally
uses the smaller acting group $I\Gamma$. Neither paper computes the
relative CM parity character of this prompt.
\end{proof}

\begin{statement}{P.6}{Open}{Search boundary for finite characters, strictness and novelty}
The opened literature search did not locate an independent theorem
proving the precise R1 lifting statement, the R2 prescribed mixed-frame
lift, the R4 relative quadratic character, or the complete R7 torsion
classification. It also did not locate a later solution of the profinite
abelianization question in \citet[Question 1.2]{LitLochakSchnepsOpenProblems}.
This is a record of the search, not a proof that no such source exists.
\end{statement}
We searched both available web indexes for the named sources, exact
titles, profinite versus pro-prime strictness, finite characters,
non-Galois gentle elements, torsion and CM comparisons. The primary
sources and exact loci are recorded in
\texttt{agent\_literature/source\_audit.md}. We checked the revision
of \citet{LitCollas2026v2} dated 5 November 2026; its strong
indecomposability assertion concerns the group rather than its
abelianization. In \citet[Proposition 5.20 and closing remarks,
p.373]{LitGuillot2014}, an unhatted inverse-limit group denotes the
usual $\GZ$; the hatted pentagon group is described as a subgroup,
without an assertion of properness. That notation cannot establish
strictness. R2 and R4 do give an actual non-Galois element of $\GZ$
under the accepted supplied results. We claim no historical priority
for that deduction or for any proof below.

\subsection{An integral construction proving strict inclusion}


\begin{statement}{P.7}{Proved}{integral peripheral realization}
\end{statement}
Let $P_\ell$ be free pro-$\ell$ on $x,y$, let
$z=(xy)^{-1}$, and let $L=\operatorname{Lie}_{\mathbf Z_\ell}(X,Y)$.
Write $\gamma_m$ for the closed ordinary lower central series. Its
associated graded Lie algebra is $L$, by the integral Magnus map.
For $n\geq2$, every homogeneous special outer derivation of degree $n$
is the leading class of a continuous multiplier-one peripheral outer
automorphism of $P_\ell$. This statement holds at every prime.

\emph{Proof.}
We first justify the two properties of the lower central comparison which
we need. The Magnus expansion sends the generators to $1+X,1+Y$;
commutators in $\gamma_m$ have leading degree $m$, and the leading
terms of the standard Lyndon commutators give the integral free-Lie
basis. Consequently the comparison is an isomorphism on every
$\gamma_m/\gamma_{m+1}$, with coefficients in $\mathbf Z_\ell$.
It is separated and complete for this filtration: every finite
$\ell$-group is nilpotent, so the intersection of the $\gamma_m$ is
trivial, and compactness identifies $P_\ell$ with
$\varprojlim_m P_\ell/\gamma_m$.

Choose a representative of the special outer derivation fixing $X$.
Its two other peripheral conjugators are homogeneous $h,k\in L_n$,
with
\[
             [h,Y]=[k,X+Y].
\]
Choose actual elements of $\gamma_nP_\ell$ with these leading terms.
Their provisional images of $y,z$, together with $x$, satisfy the
peripheral product identity through degree $n+1$. If the error is in
$\gamma_m$, alter the two conjugators in degree $m-1$. The change
in the degree-$m$ error is
\[
 L_{m-1}\oplus L_{m-1}\longrightarrow L_m,
 \qquad (u,v)\longmapsto[u,Y]+[v,-X-Y].
\]
This map is onto integrally. Its image is
$[L_{m-1},X]+[L_{m-1},Y]$, and Jacobi repeatedly moves a generator
to the outside of any bracket of total degree $m$. Thus every error
can be killed without modifying the prescribed leading terms.
The corrections converge in every finite $\ell$-group quotient,
and therefore in $P_\ell$. They give an exact peripheral endomorphism.
It induces the identity on the Frattini quotient, so the pro-$\ell$
Burnside basis theorem makes it surjective.

For completeness, the Hopficity step is as follows. If $\alpha:G\to G$
is a continuous epimorphism of a finitely generated profinite group,
then $N\mapsto\alpha^{-1}N$ injects the finite set of open normal
subgroups of a given index into itself, hence permutes it. For any
open normal $N$, some $k>0$ has $\alpha^{-k}N=N$. Since
$\ker\alpha\subseteq\ker\alpha^k\subseteq\alpha^{-k}N=N$,
the kernel lies in every open normal subgroup and is trivial.
Applying this to the endomorphism above proves the assertion.
Inner derivations are realized by inner automorphisms, so the
construction passes to outer classes. No exponential or division by
factorials is used. \hfill$\square$

\begin{statement}{P.8}{Proved}{coprime geometric symmetry lifts}
\end{statement}
For $\ell\geq5$, every integral special outer leading class fixed by
the geometric $S_3$ lifts to an actual special outer automorphism
centralizing that $S_3$.

\emph{Proof.}
Let $I^n$ be the closed subgroup of IA special outer automorphisms
whose first possible derivation degree is $n$. It is pro-$\ell$:
in each finite characteristic $\ell$-group quotient, the automorphisms
trivial on the Frattini quotient form an $\ell$-group. The geometric
$S_3$ acts by conjugation on $I^n$ and preserves $I^{n+1}$.
P.7 proves surjectivity onto the full integral special leading module.
Given a lift $a$ of an invariant class, the elements $a^{-1}g(a)$,
$g\in S_3$, form a nonabelian $1$-cocycle in $I^{n+1}$.
Its class is zero: in a finite quotient this is the conjugacy of
complements in an $\ell$-group-by-$S_3$ extension, since $6$ is
prime to $\ell$. One may prove this by induction on an invariant
central elementary abelian subgroup, averaging the additive
cocycles there by $1/6$. Compactness of the finite solution sets
gives a single correction in $I^{n+1}$ at the inverse limit.
Correcting $a$ makes it $S_3$-fixed without changing its leading
term. This is a group construction, not a claim that an arbitrary
graded solution exponentiates integrally. \hfill$\square$

\begin{statement}{P.9}{Proved}{exact normalized pro-$\ell$ pairs}
\citep[Propositions 3--4 and Theorem 2]{LitHarbaterSchneps2000}
\end{statement}
The centralizer in P.8 is identified with the multiplier-one
two-cycle and three-cycle pairs on $P_\ell$, with their actual
automorphism condition. If its first outer derivation is $-[D_h]$,
where $D_h(X)=0$ and $D_h(Y)=[h,Y]$ has degree $n\geq2$, the
normalized pair $f_\ell$ lies in $\gamma_nP_\ell$ and has leading
term exactly $h$.

\emph{Proof.}
Put $E_\ell=\widehat B/O^\ell(F)$. It is the actual relative
completion in
\[
 1\longrightarrow P_\ell\longrightarrow E_\ell
             \longrightarrow S_3\longrightarrow1,
 \qquad x=(sr)^{-2},\quad y=sxs.
\]
The center of $P_\ell$ is trivial, and $S_3$ acts faithfully on
$P_\ell^{\rm ab}$. Thus conjugation embeds $E_\ell$ into
$\operatorname{Aut}(P_\ell)$ and identifies it with the pullback
of $\operatorname{Aut}(P_\ell)\to\operatorname{Out}(P_\ell)$ along
the geometric map $S_3\to\operatorname{Out}(P_\ell)$.
An outer automorphism centralizing this map therefore extends,
uniquely up to an inner automorphism by $P_\ell$, to $E_\ell$,
inducing the identity on $S_3$.

Choose a special representative fixing $x$. Its extension sends
$\tau=sr$ to an element with the same square $x^{-1}$ and the
same image in $S_3$. Their quotient lies in $C_{P_\ell}(x)=\langle
x\rangle$. If it is $x^v$, equality of squares gives $2v=0$ in
$\mathbf Z_\ell$, hence $v=0$. Thus $\tau$ is fixed literally.
The centralizer assertion follows from the Magnus leading terms:
after removing the abelian power of $x$, a first nonzero term
commuting with $X$ would belong to the centralizer of $X$ in
$L_m$, which is zero for $m>1$. That last assertion follows by
embedding in the free associative algebra; a homogeneous tensor
commuting with $X$ is a multiple of $X^m$, and no positive
commutator Lie element has such a term.

By coprime complement conjugacy, the new $s,r$ are conjugate to
the old positive generators by elements of $P_\ell$. Put $f=s\Phi(s)$.
Since $\Phi(s)\Phi(r)=sr$, we have
\[
 \Phi(s)=sf,\qquad \Phi(r)=f^{-1}r,
 \qquad \Phi(y)=f^{-1}yf.
\]
The first torsion relation makes $f^{\rm ab}=(a,-a)$. Replacing
$\Phi$ by $\operatorname{Inn}(x^a)\Phi$ keeps $\tau$ fixed and
replaces $f$ by $y^af x^{-a}\in P_\ell'$. In the marked group
conjugation by $s$ exchanges $x,y$, and conjugation by $r$
cycles $x,y,z$, where $z=(xy)^{-1}$. The two torsion relations
therefore give, with this exact order,
\[
 f(y,x)f(x,y)=1,\qquad f(z,x)f(y,z)f(x,y)=1.
\]
The restriction of $\Phi$ is already an automorphism. Conversely,
these equations and that automorphism condition define an
automorphism of $E_\ell$ by the displayed assignments. Its outer
restriction centralizes $S_3$ because the quotient action is the
identity. This also gives a direct proof of the pro-$\ell$ version
of Harbater--Schneps, Propositions 3--4 and Theorem 2, rather than
silently transporting their profinite result to a different category.

It remains to verify the leading coefficient. In
$P_\ell/\gamma_{n+1}$, the centralizer of $x$ is
$\langle x\rangle\gamma_n/\gamma_{n+1}$, by the same leading-term
argument. If $\psi_f$ is inner there, say by $w$, then $w$ belongs
to that centralizer and $fw$ belongs to the corresponding centralizer
of $y$. Since $f^{\rm ab}=0$, both powers vanish and
$f\in\gamma_n$. If its leading term is $h_f$, then
$\psi_f(y)y^{-1}$ has leading term $-[h_f,Y]$.
Equality of the prescribed outer class with $-[D_h]$, followed by
the free-Lie centralizer argument at $X$ and then at $Y$, gives
$h_f=h$. All maps here are continuous quotient or conjugation maps.
\hfill$\square$

\begin{statement}{P.10}{Proved}{an independently checked nonpentagonal leading term}
\citep[equations (1.2)--(1.4)]{LiftCohen2001}
\end{statement}
For the standard Lyndon brackets $b_w$ with $x<y$, set
\[
\begin{split}
q_7={}&3b_{xxxyxyy}+3b_{xxxyyxy}+b_{xxyxyxy}
       -3b_{xxyxyyy}\\
     &-3b_{xxyyxyy}-3b_{xxyyyxy}-b_{xyxyxyy}.
\end{split}
\]
It has integral coefficients, is nonzero, and satisfies
\[
 q_7(x,y)+q_7(y,x)=0,\qquad
 [x,q_7(x,y)]=[x+y,q_7(-x-y,y)].
\]
It is special and has zero three-cycle norm. Its infinitesimal
pentagon defect is nonzero, already over every field of
characteristic at least five.

\emph{Proof and exact certificate.}
The program \texttt{check\_lie\_witnesses.py} expands standard
Lyndon brackets by splitting at the longest proper Lyndon suffix.
All four asserted identities are checked coefficientwise over the
integers; $q_7$ has 62 nonzero associative coefficients. Here are the
inner terms needed for the symmetry assertion. Put
$\alpha(X)=Y$, $\alpha(Y)=X$, and
$\theta(X)=-X-Y$, $\theta(Y)=Y$. They are the transpositions
$(X,Y)$ and $(X,Z)$ on $X+Y+Z=0$ and generate $S_3$. If
$D_h(X)=0$, $D_h(Y)=[h,Y]$ and $k=\theta(h)$, antisymmetry and
the displayed cabling equation give exactly
\[
 \alpha D_h\alpha^{-1}-D_h=\operatorname{ad}(-h),\qquad
 \theta D_h\theta^{-1}-D_h=\operatorname{ad}(k-h).
\]
For example the latter difference sends $X$ to $-[k,Y]$ and
$Y$ to $[k-h,Y]$, while cabling says
$[k-h,X]=-[k,Y]$. Thus $[D_h]$ is fixed by both generators.
The program additionally checks these two derivation identities on
both free generators for $q_7$ and $h_9$.

We retain an actual braid comparison. The kernel of
$P_4\to P_3$ is free on the three braids $A_{14},A_{24},A_{34}$.
Conjugation by $P_4$ on this kernel, followed by its pro-$\ell$
completion, is an actual continuous group action; the full twist
acts trivially, so it factors through the pure five-point sphere
group. Its leading derivations on $L(a,b,c)$ are
\[
 t_{12}\mapsto D_{ab},\quad t_{13}\mapsto D_{ac},\quad
 t_{23}\mapsto D_{bc},\qquad
 t_{14}\mapsto\operatorname{ad}a,\quad
 t_{24}\mapsto\operatorname{ad}b,\quad
 t_{34}\mapsto\operatorname{ad}c,
\]
where $D_{uv}(u)=[u,v]$, $D_{uv}(v)=[v,u]$, and $D_{uv}$
kills the third generator. These formulas follow by linearizing
the conjugation presentation in Cohen, equations (1.2)--(1.4).
They are therefore leading terms of the actual group action,
not merely a representation of an unrelated presented Lie algebra.
The program independently checks all 18 required infinitesimal
braid relations on these derivations.

For the pentagon convention
\[
\begin{split}
\mathcal P(h)={}&h(t_{12},t_{23}+t_{24})
 +h(t_{13}+t_{23},t_{34})-h(t_{23},t_{34})\\
 &-h(t_{12}+t_{13},t_{24}+t_{34})-h(t_{12},t_{23}),
\end{split}
\]
the program obtains 880, 882 and 808 nonzero coefficients on
$a,b,c$, respectively. The coefficient of $a^5bcb$ in
$\mathcal P(q_7)(a)$ is exactly $16$. The complete 2,570-coefficient
output is retained in \texttt{lie\_witnesses\_exact.json}.
This coefficient proves nonvanishing over every $\mathbf Q_\ell$
and modulo every $\ell\geq5$. \hfill$\square$

\begin{statement}{P.11}{Proved}{strictness for the full profinite groups}
\end{statement}
There is an element of $\widehat{\mathrm{GT}}_0$ outside
$\widehat{\mathrm{GT}}$. In particular
\[
             \widehat{\mathrm{GT}}
                    \subsetneq\widehat{\mathrm{GT}}_0.
\]
This proves the corollary requested in P.2. It assigns no value to
$\Delta$ on $\widehat{\mathrm{GT}}$.

\emph{Proof.}
Fix any $\ell\geq5$. P.7--P.9 lift the invariant class
$-[D_{q_7}]$ to an actual normalized pair
$f_\ell\in\gamma_7P_\ell\cap P_\ell'$ whose leading term is $q_7$.
Its pure map is an automorphism and it satisfies both group cycle
relations. The last assertion of the refereed R1 applies to
precisely this pair, retaining its entire pro-$\ell$ image. It
gives $(1,f)\in\widehat{\mathrm{GT}}_0$ with
$q_\ell(f)=f_\ell$.

If this full pair satisfied the pentagon, apply the continuous
quotient to pro-$\ell$ pure braid groups and then the action on
the free pro-$\ell$ kernel used in P.10. Taking the first
nonzero Johnson degree of the resulting identity gives
$\mathcal P(q_7)=0$. To justify this direction explicitly,
substitution of a group element whose first term has degree seven
depends in that degree only on the degree-one terms of its
arguments. Products of the five substituted elements add in that
degree; all commutators between them and all higher corrections
have larger degree. The cabling and quotient homomorphisms used
in the group pentagon induce exactly the sums of $t_{ij}$ in
P.10. Thus a zero group defect forces a zero displayed leading
derivation. Its coefficient 16 contradicts that conclusion.
Every full lift of this same pro-$\ell$ pair is excluded for this
reason; no inference from a surviving finite shadow is involved.
\hfill$\square$

\subsection{The relative abelianization}


\begin{statement}{P.12}{Proved}{the precise abelianization question}
\citep[Question 1.2]{LitLochakSchnepsOpenProblems}
\end{statement}
For either $G=\widehat{\mathrm{GT}}$ or
$G=\widehat{\mathrm{GT}}_0$ let
$G^{\rm ab}=G/\overline{[G,G]}$ and
$K_G=\ker(G^{\rm ab}\xrightarrow{\lambda}\widehat{\mathbf Z}^{\times})$.
Kronecker--Weber identifies
$G_{\mathbf Q}^{\rm ab}$ with $\widehat{\mathbf Z}^{\times}$
via $\chi$.  The inclusion $G_{\mathbf Q}\hookrightarrow G$
therefore gives a continuous section of the displayed multiplier
map.  Thus
\[
 G^{\rm ab}\simeq\widehat{\mathbf Z}^{\times}\times K_G
\]
as profinite abelian groups, with the Galois abelianization as the
first factor.  Every continuous finite abelian character of
$G_{\mathbf Q}$ extends to $G$ by composition with $\lambda$.
Characters of $G$ trivial on $G_{\mathbf Q}$ are exactly the
continuous characters of $K_G$.  For finite-order complex
characters this is its Pontryagin dual.

Lochak--Schneps, \emph{Open problems in Grothendieck--Teichm\"uller
theory}, Question 1.2, asks whether $K_{\widehat{\mathrm{GT}}}=1$.
Equivalently, whether every finite-order character of
$\widehat{\mathrm{GT}}$ factors through $\lambda$: characters
separate points of a profinite abelian group, since every nonidentity
element survives in a finite abelian quotient and in a character
of that quotient.  R4 proves that
$K_{\widehat{\mathrm{GT}}_0}$ has a quotient $C_2$.
This last fact neither computes that kernel nor decides whether
the restriction of this quotient to $\widehat{\mathrm{GT}}$
is nonzero.  We found no classification of the finite characters
or resolution of Question 1.2 in the primary sources searched.

\subsection{A strict inclusion inside the sign kernel}


\begin{statement}{P.13}{Proved}{The kernel of the sign is strictly larger than its pentagon subgroup}
Under the supplied R1, R2 and R4, both fibres of
$\Delta:\GZ\to\Ftwo$ contain actual multiplier-one elements
which fail the profinite pentagon. In particular,
\[
 G_{\Q}\ \subsetneq\ \ker\Delta,
 \qquad
 \GT\cap\ker\Delta\ \subsetneq\ \ker\Delta.
\]
If the main vanishing assertion holds, its inclusion chain is
therefore necessarily
\[
 G_{\Q}\ \subseteq\ \GT\ \subsetneq\ \ker\Delta
                       \ \subsetneq\ \GZ.
\]
Only the first inclusion in that conditional chain remains of
unknown strictness here. This statement does not decide the
main vanishing assertion.
\end{statement}
\begin{proof}
Fix $\ell\ge5$ and let $g\in\HH$ be the actual R1 lift in
P.11 of the independently constructed nonpentagonal
pro-$\ell$ pair. Let $k\in\HH$ be the R2 element, with
$\Delta(k)=1$ and identity action on
$\widehat B/R_{\rm nil}(U)$, where $U=F'F^3$.
For every homomorphism from $F$ to a finite $\ell$-group,
its restriction to $U$ has nilpotent image. Thus
\[
 R_{\rm nil}(U)\subseteq O^\ell(F).
\]
The action of $k$ on the actual relative pro-$\ell$ quotient
$\widehat B/O^\ell(F)$ is consequently identity.
Its normalized pro-$\ell$ parameter is one: it centralizes
$y$, belongs to $P_\ell'$, and
$C_{P_\ell}(y)=\langle y\rangle$ has trivial intersection
with that derived subgroup. Hence multiplying $g$ on either
side by any power of $k$ retains the complete normalized
pro-$\ell$ pair, not just its leading term.

Write $e=\Delta(g)\in\{0,1\}$ and set
$h_0=gk^e$, $h_1=h_0k$. Both are actual elements of $\HH$.
The character law gives $\Delta(h_i)=i$, while their
pro-$\ell$ pentagon defects retain the nonzero coefficient
$16$ from P.10. A profinite pentagon solution maps to a
pro-$\ell$ pentagon solution by the actual braid quotient
used there, so neither $h_i$ is in $\GT$.
In particular $h_0\in\ker\Delta\setminus\GT$ and is not
Galois. This proves both strict inclusions. If
$\Delta|_{\GT}=0$, replace $\GT\cap\ker\Delta$ by $\GT$;
the final strict inclusion follows from $\Delta(k)=1$.
The proof needs no guessed value of $\Delta(g)$: the two
case-dependent definitions select the required element
using the already defined finite continuous character.
\end{proof}
